\chapter{拖动示教 回放与运动学习}
\label{ch:teaching}

示教的本质是把人的意图转换成一条可以复现的参考轨迹。拖动时的力矩控制、数据录制、轨迹修整和回放时的跟踪控制，是四个不同问题。把它们拆开，才容易判断“推起来不顺”“回放不平滑”和“回放位置不准”分别出在哪里。

\section{拖动时机器人到底应该做什么}

对力矩可控机械臂，最简单的教学模型是
\begin{equation}
\tau_{\rm cmd}=\hat g(q)+\hat\tau_f(\dot q)-D_q\dot q.
\end{equation}
$\hat g$ 估计并补偿重力；$\hat\tau_f$ 补偿部分摩擦；$D_q\succeq0$ 提供阻尼。理想补偿且没有外力时，剩余运动方程为 $M\ddot q+C\dot q+D_q\dot q=0$。半正定阻尼只保证不增加机械能，未阻尼的运动模态可能持续存在。若 $D_q\succ0$，并满足轨迹有界等常规条件，才可用能量与不变集分析得到速度趋于零；系统没有位置弹簧，不会自动回到某个固定构型。

\textbf{阻尼不能在静止时托住重力。}因为 $\dot q=0$ 时 $-D_q\dot q=0$。如果重力补偿少了 $1\,\mathrm{N\,m}$，加一点阻尼只能限制下落速度，不能使任意姿态静止平衡。以一维近似 $J\ddot q+D\dot q=-1$ 为例，稳态速度为 $-1/D$。取 $D=2\,\mathrm{N\,m\,s/rad}$，仍会以 $-0.5\,\mathrm{rad/s}$ 漂移。原稿中把小阻尼描述成“托住”的说法需要这样修正。

若采用导纳拖动，先由测得或估计的外力生成参考位移，再交给位置内环：
\begin{equation}
M_v\ddot x_r+D_v\dot x_r+K_v(x_r-x_0)=F_{\rm ext}.
\end{equation}
$K_v=0$ 时松手后速度会衰减，但停留位置由此前运动决定；$K_v>0$ 时会回到 $x_0$。关节拖动、末端某几个方向自由、肘部自由，可以用不同任务约束实现，不能把一个标量“软硬”旋钮当成全部设计。

\section{记录哪些量才能在以后解释结果}

每条样本至少包括单调时间戳、按固定关节名排序的实测 $q$、控制模式以及数据有效标记。建议同时记录实测或估计速度、指令力矩、电流换算力矩、末端位姿和所用模型版本。关节命令不能替代关节测量；电流换算力矩也不能在没有校准时当成独立高精度力矩传感器。

\begin{table}[htbp]\centering\small
\begin{tabularx}{\linewidth}{p{31mm}XX}\toprule
字段 & 用途 & 常见错误\\\midrule
$t_k$ 与序号 & 重建真实时间和丢帧 & 用数组下标假定每帧等间隔\\
关节名与 $q_k$ & 保留实际构型 & 左右臂或关节顺序串位\\
参考与测量 & 分开路径问题和跟踪问题 & 只保存指令画出“完美跟踪”\\
工具与基座变换 & 重建实际任务坐标 & 换工具后继续用旧 TCP\\
质量与质心版本 & 解释重力和观测器残差 & 只记录文件名没有参数版本\\
模式与限幅状态 & 找切换和饱和问题 & 把被截断的输出当控制器原输出\\\bottomrule
\end{tabularx}
\caption{一条轨迹文件应该同时保存轨迹本身及解释它所需的上下文。}
\end{table}

假定录制名义频率为 $100\,\mathrm{Hz}$，真实样本时间是 $0,0.010,0.023\,\mathrm{s}$，对应位置 $0,0.002,0.006\,\mathrm{rad}$。在 $0.020\,\mathrm{s}$ 重采样，线性插值为
\[
q=0.002+\frac{0.020-0.010}{0.023-0.010}(0.006-0.002)
=0.0050769\,\mathrm{rad}.
\]
若误把第三帧当成 $0.020\,\mathrm{s}$，就直接取 $0.006$，人为引入约 $0.0009231\,\mathrm{rad}$ 差异。以后再对它求速度，时间错误还会继续放大。

\section{把一百赫兹记录交给两百赫兹控制器}

插补到 $200\,\mathrm{Hz}$ 不会创造新的测量信息，只是在已有样本之间构造参考。步骤是：去掉重复或倒退时间戳，按关节名映射，检查跳变，选择拟合或滤波方法，再在控制器自己的时钟上查询轨迹。它可以提高参考的连续性，但不能把原始 $100\,\mathrm{Hz}$ 信号的可恢复频率变成 $100\,\mathrm{Hz}$。

\subsection{平滑和保形不是同一个要求}

三次样条通常具有连续的二阶导数，但可能在稀疏点之间过冲。PCHIP 是分段三次 Hermite 保形插值，在标量单调数据区间可以保持单调性，通常只有一阶导数连续。把它说成“同时没有过冲且加速度处处连续”是不正确的。Akima 也是局部三次方法，但不能一般性保证单调和无过冲。

若关节必须满足限位，应该检查连续插值曲线的极值；只检查记录点不够。如果必须让加速度连续，可采用带边界条件的平滑样条或优化求解，并把限位作为约束。对姿态则需要在旋转空间插值；先统一四元数符号，再进行 SLERP 或适当的姿态样条。

\begin{figure}[htbp]\centering
\includegraphics[width=0.94\linewidth]{figures/rt_teach.pdf}
\caption{示教轨迹的采样、噪声和重建示例。左图是位置，右图比较差分速度与平滑后速度。平滑会改变原数据，必须同时看位置偏差与动态峰值，不能只看曲线是否好看。}
\end{figure}

离线双向滤波可以避免净相位滞后，但必须使用未来数据，并会改变幅值与边界。处理前后各保留一份数据，裁边范围和滤波参数写入轨迹元数据。如果要保留关键时刻的接触动作，不应盲目把所有高频成分都删掉。

\subsection{回放起点和时间缩放}

回放前，当前位置 $q_{
m now}$ 通常不同于轨迹第一点 $q_0$。需要先生成一段从 $q_{
m now}$ 到 $q_0$ 的过渡轨迹，完成后再启用录制时间轴。把第一帧直接发给内环，会产生一个不在示教数据里的跳变。

若把两秒示教压到一秒，位置路径不变，速度加倍、加速度四倍、jerk 八倍。假定原峰值分别为 $0.4\,\mathrm{rad/s}$、$0.8\,\mathrm{rad/s^2}$、$2\,\mathrm{rad/s^3}$，压缩后就是 $0.8,3.2,16$。这也是为什么“人手推得动”不能作为自动高速回放的可行性证明。

\section{关节回放和笛卡尔回放如何选择}

关节回放直接复用 $q_d(t)$，保留示教时的肘部构型，避免重新选择逆解分支。但实际轨迹仍可能因控制误差、弹性、障碍变化和工具变化而偏离；它不保证绝对定位精度，也不自动保证无碰撞。

笛卡尔回放保存 $T_{BE}(t)$，方便把整个动作变换到新的工件坐标系。例如原工作坐标系是 $W$，工具轨迹为 $T_{WE}(t)$；新工作坐标系在基座中的位置是 $T_{BW}'$，那么新目标是
\begin{equation}
T_{BE}'(t)=T_{BW}'T_{WE}(t).
\end{equation}
这一步是坐标变换，不是把每个坐标数直接加上一个偏移。之后仍要做连续逆解、冗余姿态处理和完整可行性检查。

\begin{examplebox}[把运镜轨迹绕基座旋转]
原轨迹一点在工作系中为 $p_W=(0.2,0.1,0)\,\mathrm m$。新工作系相对基座绕 $z$ 旋转 $90^\circ$，原点为 $(0.5,0,0)$。于是
\[
p_B=\begin{bmatrix}0&-1&0\\1&0&0\\0&0&1\end{bmatrix}
\begin{bmatrix}0.2\\0.1\\0\end{bmatrix}
+\begin{bmatrix}0.5\\0\\0\end{bmatrix}
=\begin{bmatrix}0.4\\0.2\\0\end{bmatrix}\mathrm m.
\]
如果只把工作系原点加上去，就会错误地得到 $(0.7,0.1,0)$。工具姿态也必须左乘同一个旋转矩阵。
\end{examplebox}

\section{动态运动基元从稳定系统加形状项得到}

\textbf{DMP（dynamic movement primitive，动态运动基元）}把示教形状写进一个动力系统。这里用 $y$ 表示一个标量运动坐标，$g$ 表示目标，$\rho$ 表示从一衰减到零的相位，$T_d>0$ 是时间尺度，取 $\alpha_z>0,\beta_z>0,\alpha_\rho>0,\rho(0)=1$：
\begin{align}
T_d\dot z&=\alpha_z\big(\beta_z(g-y)-z\big)+f(\rho),\\
T_d\dot y&=z,\\
T_d\dot\rho&=-\alpha_\rho\rho.
\end{align}
消去 $z$ 得
\begin{equation}
T_d^2\ddot y+\alpha_zT_d\dot y+\alpha_z\beta_z(y-g)=f(\rho).
\end{equation}
当 $f=0$，它就是一个稳定的弹簧阻尼系统；取 $\beta_z=\alpha_z/4$ 时，对应临界阻尼。相位解为 $\rho(t)=\exp(-\alpha_\rho t/T_d)$。形状项采用高斯基函数
\begin{equation}
\psi_i(\rho)=\exp[-h_i(\rho-c_i)^2],\qquad
f(\rho)=\frac{\sum_i\psi_i(\rho)w_i}{\sum_i\psi_i(\rho)}\rho a_s.
\end{equation}
$c_i$ 是中心，$h_i>0$ 控制宽度，$w_i$ 是要学习的权重，$a_s$ 是空间尺度。只要权重有界、分母有效、目标固定，且没有破坏稳定性的额外耦合，$\rho\to0$ 会让形状项消失，最终回到稳定目标系统。加入避障耦合、时变目标或真实执行器之后，需要重新分析，不能再用一句“DMP一定稳定”代替证明。

\subsection{权重怎样从示教数据算出来}

已有示教 $y_d,\dot y_d,\ddot y_d$ 时，先把它代入方程反求所需形状力：
\begin{equation}
f_d=T_d^2\ddot y_d-\alpha_z\big(\beta_z(g-y_d)-T_d\dot y_d\big).
\end{equation}
定义 $\Phi_{ki}=\rho_k a_s\psi_i(\rho_k)/\sum_j\psi_j(\rho_k)$，于是 $f_d\approx\Phi w$，求解 $\min_w\|\Phi w-f_d\|^2+\lambda\|w\|^2$。这正是前面学过的正则化最小二乘，DMP 没有绕开辨识数学。

\begin{examplebox}[一个 DMP 样本怎样产生回归目标]
取 $T_d=1,\alpha_z=4,\beta_z=1,g=1$。某时刻示教为 $y_d=0.5,\dot y_d=1,\ddot y_d=0$，则
\[
f_d=0-4[(1-0.5)-1]=2.
\]
若只用一个基函数，此处 $\rho=0.5,a_s=1$，归一化后 $\Phi=0.5$，因此这个样本要求 $0.5w=2$，即 $w=4$。多个时刻通常不能被同一个权重精确满足，所以要堆叠成回归问题，而不是每个时刻单独改一次权重。
\end{examplebox}

常见写法令 $a_s=g-y_0$。如果任务起点等于终点，尺度变成零，整个形状项会丢失。可以使用独立尺度、分段动作或周期型 DMP，不能只在除法里加一个很小的数就声称问题已解决。

\begin{figure}[htbp]\centering
\includegraphics[width=0.94\linewidth]{figures/rt_dmp.pdf}
\caption{同一组 DMP 权重的重建与时间缩放示例。改变时间尺度会改变速度；改变目标或尺度后仍应重新检查关节和碰撞限制。曲线由方程数值积分得到。}
\end{figure}

\section{概率运动基元与一次条件化计算}

\textbf{ProMP（probabilistic movement primitive，概率运动基元）}不只学一条曲线，而是学权重分布。设
\begin{equation}
y(t)=\phi(t)^\mathsf Tw,\qquad w\sim\mathcal N(\mu,\Sigma).
\end{equation}
那么 $y(t)$ 的均值为 $\phi^\mathsf T\mu$，方差为 $\phi^\mathsf T\Sigma\phi$。不同示教在某段很一致，通常使该段方差较小；但这也可能只是示教样本不足，不能直接把方差当成任务重要程度的唯一证据。

要让时刻 $t_*$ 经过观测值 $y_*$，并假设观测噪声方差为 $r$，高斯条件公式为
\begin{align}
K&=\frac{\Sigma\phi_*}{\phi_*^\mathsf T\Sigma\phi_*+r},\\
\mu'&=\mu+K(y_*-\phi_*^\mathsf T\mu),\\
\Sigma'&=\Sigma-K\phi_*^\mathsf T\Sigma.
\end{align}
它与卡尔曼滤波的测量更新形式相同，但这里估计的是轨迹权重。

取 $\phi(t)=(1,t)^\mathsf T,\mu=(0,1)^\mathsf T,\Sigma=\operatorname{diag}(0.04,0.09)$，在 $t_*=0.5$ 希望经过 $0.8$，令 $r=0.01$。原预测均值为 $0.5$、方差为 $0.0625$，因此
\[
K=\frac{(0.04,0.045)^\mathsf T}{0.0725}
=(0.551724,0.620690)^\mathsf T,
\]
\[
\mu'=(0.165517,1.186207)^\mathsf T.
\]
新轨迹在 $0.5$ 时的均值为 $0.758621$，没有精确到 $0.8$，因为我们允许观测有噪声。取 $r\to0$ 才趋向硬约束，但多条相互矛盾的硬约束可能无法同时满足。

\section{动态时间规整与重复学习}

\subsection{DTW 先对齐动作进度}

\textbf{DTW（dynamic time warping，动态时间规整）}寻找两条序列的单调对应关系。给 $a_i,b_j$，局部代价 $d_{ij}=(a_i-b_j)^2$，累计代价满足
\begin{equation}
D_{ij}=d_{ij}+\min(D_{i-1,j},D_{i,j-1},D_{i-1,j-1}).
\end{equation}
边界设 $D_{00}=0$，其它零行零列为正无穷，再从 $(1,1)$ 计算。对 $a=(0,1,2)$、$b=(0,0,1,2)$，可以将第一个零与两个零对应，总代价为零。按数组下标硬比则会把速度差误判成形状差。实际应限制匹配带宽，防止把明显错误动作过度拉伸成“匹配”；DTW 代价也不是以秒定义的时序跟踪误差。

\subsection{ILC 在重复任务之间修正前馈}

\textbf{迭代学习控制（iterative learning control，ILC）}使用第 $r$ 次完整执行的误差，更新第 $r+1$ 次的输入。对离散整段任务写 $y_r=Pu_r+d$，其中 $P$ 是输入到输出的线性映射，$d$ 是每次相同的扰动。令 $e_r=y_d-y_r$，用
\begin{equation}
u_{r+1}=u_r+Le_r,
\end{equation}
则
\begin{equation}
e_{r+1}=(I-PL)e_r.
\end{equation}
标量例中 $P=0.8,L=1$，误差每次乘以 $0.2$。若首次为 $0.1\,\mathrm{rad}$，后两次为 $0.02$ 和 $0.004\,\mathrm{rad}$。但如果 $L=3$，乘子为 $-1.4$，误差会变号并放大。ILC 是跨次学习，不是实时积分；初始条件和接触情况变化过大时，上一轮学习可能失效。

\section{怎样评价回放}

同一时间轴上的轨迹跟踪误差可用
\begin{equation}
E_{q,\rm RMS}=\sqrt{\frac1N\sum_{k=1}^N\|q_k-q_{d,k}\|^2},\quad
E_{p,\rm RMS}=\sqrt{\frac1N\sum_{k=1}^N\|p_k-p_{d,k}\|^2}.
\end{equation}
不同单位不要混进同一个未加权范数。\textbf{重复性}描述多次到同一目标时的离散程度；\textbf{准确度}描述均值与真实目标之间的偏差。每次都偏 $5\,\mathrm{mm}$、散布只有 $0.2\,\mathrm{mm}$，代表重复性好但存在系统偏差。仅用同一个错误运动学模型计算参考和实际 FK，发现不了真实末端的安装和标定误差。

\section{计算练习与解答}

\subsection*{题一 时间偏差折成位置偏差}
某关节近似匀速 $0.2\,\mathrm{rad/s}$，参考和测量错位 $10\,\mathrm{ms}$。求伪误差及 $0.5\,\mathrm m$ 力臂上的近似位移。

解：$\delta q=0.2\times0.01=0.002\,\mathrm{rad}$，小角度位移 $\delta p\approx0.5\times0.002=0.001\,\mathrm m=1\,\mathrm{mm}$。毫米级比较先要求毫秒级时间对齐。

\subsection*{题二 DMP 的时间尺度}
保持同一相位轨迹，把 $T_d$ 从 $2$ 改为 $4\,\mathrm s$。原峰值速度为 $0.8$，加速度为 $1.6$，求新值。

解：整体时间增加两倍，速度为 $0.4$，加速度为 $0.4$。对实际机器人还应检查是否在新速度区间进入更明显的静摩擦或爬行状态。

\subsection*{题三 ILC 的收敛范围}
标量对象 $P=0.8$，求使误差收敛的学习增益 $L$。

解：要求 $|1-0.8L|<1$。展开得 $-1<1-0.8L<1$，因此 $0<L<2.5$。端点 $L=0$ 不学习，$L=2.5$ 误差等幅变号，均不渐近收敛。

\par\smallskip
{\small\noindent\textbf{来源与扩展}\quad 本章对应 [B，第93--97页]，承接 [A，第52--60页]；补充示教时序、插值、DMP/ProMP与ILC计算，并明确回放精度与重力补偿的条件。\par}
